% TC-01 Sensor tow stability, blower rig. Rev B replaces the rev A car tow.
% Build: ../build.sh
\def\dbfroot{../../design-system}
\documentclass[final]{dbfdoc}
\dbfsetdoc{Sensor Tow Stability}{Test Card TC-01}{AUM-DBF-27-TC-001}{C}
\dbfsetdate{11 Oct 2026}
\raggedbottom

\begin{document}

\begin{frontmatter}
\section{Objective}
\noindent\begin{minipage}[t]{0.62\textwidth}
\setlength\parindent{1.15em}\noindent
Screen 2--3 candidate sensor bodies, each at three ballast states, for aerodynamic stability on a tow line: no spinning, rotation or sustained oscillation. Measure the line angle to estimate sensor drag for the tow model (Script~5). Film every run for the proposal.

\noindent\textbf{Rig.} The university has no wind tunnel, and the rev~A car tow is withdrawn for safety. This revision uses a leaf blower's jet, straightened in a short duct. Speed is read with an anemometer, so the same card works in a makeshift tunnel.

\noindent\textbf{Limitation.} The jet is slower and smaller than flight conditions, and a short line is not the in-flight pendulum. Stability is a screen and drag is semi-quantitative; the towed-flight test qualifies the sensor. A fin misalignment makes the body roll at a steady rate that rises in proportion to airspeed, so any rotation here would be faster in flight.
\end{minipage}\hfill
\begin{minipage}[t]{0.35\textwidth}
\vspace{-11pt}
\begin{dbfpanel}[Test at a Glance]
\dbfreadout{Crew}{3}{people}\hfill\dbfreadout{Duration}{3--4}{h}\dbfreadoutgap
\dbfreadout{Bodies}{2--3}{}\hfill\dbfreadout{Ballast states}{3}{each}\dbfreadoutgap
\dbfreadout{Duct}{100--150}{mm}\hfill\dbfreadout{Site}{Indoor}{}
\end{dbfpanel}
\end{minipage}
\end{frontmatter}

\section{Rules That Apply}
\begin{itemize}
\item §3.1.1: the sensor must be ``aerodynamically stable during deployment, operation and recovery'', which the rules define as ``remaining in a fixed orientation in the horizontal plane (not spinning or rotating) with no effect on the stability of the aircraft, at the discretion of the Flight Director.''
\item §3.1.1: minimum length 152.4~mm (6~in). Fins stay fixed (not removed for packing, not deployed), the external shape is the same for all missions, and ballast is internal.
\end{itemize}

\section{Test Articles and Equipment}
\begin{itemize}
\item \textbf{Sensor bodies:} 2--3 fixed-geometry shapes from Sensor Systems (for example a ring tail, four fixed fins, or a fixed drag cone), at least 170~mm long for margin over the 152.4~mm minimum, each with an internal ballast bay. Paint a \textbf{contrasting longitudinal stripe} on each so any roll shows on video. Mark the tow point and CG.
\item \textbf{Ballast:} light, mid and heavy internal states spanning the analysis range, weighed to 0.1~g on the bench scale (A01).
\item \textbf{Blower} (PR-01 A07), clamped down, feeding a rigid 100 or 150~mm duct, 0.6--1.0~m long (A08). Drinking straws (A09) packed tight across the duct near its exit form the straightener.
\item \textbf{Mount and line:} a rigid boom above the jet with an adjustable tie-off. Braided PE line (A11), 0.3--0.6~m long, runs to a ball-bearing swivel (A10) at the tow point, plus a secondary tether. The swivel lets a roll tendency show as rotation instead of being hidden by line twist.
\item \textbf{Measurement:} vane anemometer (A03); a printed 10~mm grid behind the line, square to the side camera; two phones at 60~fps, side-on (line angle, pitch) and from above (roll stripe, yaw); air temperature and pressure (a phone app is enough); optionally the hanging scale (A02) at the tie-off for tension.
\end{itemize}

\begin{table}[h]
\caption{Expected mean exit speed $V = Q/A$ (m/s) before straightener and duct losses. A smaller duct is faster but its uniform core is narrower; the jet survey gives the real values.}\label{tab:jet}
\begin{dbftable}{colspec={X[l] Q[r] Q[r] Q[r]}}
Air volume $Q$ & Ø100 mm & Ø125 mm & Ø150 mm \\
10 m³/min & 21.2 & 13.6 & 9.4 \\
15 m³/min & 31.8 & 20.4 & 14.1 \\
20 m³/min & 42.4 & 27.2 & 18.9 \\
\end{dbftable}
\end{table}

\section{Safety}
\begin{itemize}
\item Eye protection within 3~m of the jet; hearing protection beside the blower (A05, A06).
\item Nobody stands downstream. Keep 5~m behind the body clear: a body that breaks free leaves at jet speed.
\item Check the tether, knots, swivel and tie-off before each run. Stop the blower before touching anything in the jet.
\item Route the cord away from the jet. Run indoors or sheltered, because outdoor wind disturbs the jet.
\end{itemize}

\section{Procedure}
\begin{enumerate}
\item Photograph, weigh and measure each body and ballast state (length, diameter, fin span, mass, tow point, CG) and log them.
\item \textbf{Jet survey}, no body fitted. At the body station, about one duct diameter past the exit, traverse the anemometer across the jet horizontally and vertically in 15--20~mm steps. The \textbf{core} is where the speed is within 5\,\% of the centre value. The vane averages over its own diameter, so the core edge is approximate. Use Table~\ref{tab:jet} to pick the duct, and repeat for each blower setting.
\item Hang the body and adjust the tie-off until the whole body, fins included, sits inside the core at speed.
\item For each body, ballast state and setting, hold 30~s steady with both cameras running. Classify the run: \textbf{S} stable, \textbf{O} oscillating (amplitude, period), \textbf{C} coning, \textbf{R} rotating (fail). Log roll as none, intermittent or continuous.
\item At the top setting, flick the line once above the jet with a stick, hands outside the jet. Time the decay from the video.
\item Read the line angle $\theta$ from horizontal in three or more side-camera frames, plus the tension $T$ if fitted. Log air temperature and pressure.
\item Repeat one run of the best body with the fixed attachment planned for flight in place of the swivel.
\end{enumerate}

\section{Data Reduction}
With the line $\theta$ above horizontal, tension $T$, weight $W$ and lift $L$:
\begin{equation}
D = T\cos\theta, \qquad W - L = T\sin\theta, \qquad C_D A = \frac{2D}{\rho V^2} .
\end{equation}
Without a tension reading, take $L \approx 0$, so $D \approx W\cot\theta$. Air density is $\rho = p/(R\,T_\mathrm{air})$ with $R = 287.05$~J/(kg\,K), $p$ in Pa (1~hPa = 100~Pa) and $T_\mathrm{air}$ in K. A near-vertical line amplifies angle errors, since $\Delta D/D = \Delta\theta/(\sin\theta\cos\theta)$: a 1° misreading moves $D$ about 10\,\% at 80° and 20\,\% at 85°. Average the frames.

\section{Pass Criteria and Use of Results}
\textbf{Pass:} no rotation (R) at any setting or ballast state, and oscillation after the flick dies out within about 2~s. Passing bodies go to the towed-flight test.

\textbf{Use:} the shape, tow point and ballast range define SR-08 (draft Table~04). The measured $C_D A$ replaces the assumed sensor drag in Script~5, and the best run is the proposal photo. File the sheet as \texttt{testing/data/TC-01\_<date>.csv} and the media in \texttt{testing/media/TC-01/}.

\twocolumn[{%
\section{Data Sheet}
\noindent{\small Date \rule{1.2in}{0.4pt}\hfill Site \rule{2.4in}{0.4pt}\hfill Crew \rule{2.3in}{0.4pt}}\par\vspace{7pt}
\noindent{\small Blower and setting \rule{1.9in}{0.4pt}\hfill Duct Ø (mm) \rule{0.7in}{0.4pt}\hfill Straightener length (mm) \rule{0.8in}{0.4pt}}\par\vspace{7pt}
\noindent{\small Air temperature \rule{0.8in}{0.4pt}\hfill Pressure (hPa) \rule{0.8in}{0.4pt}\hfill Video file IDs \rule{2.7in}{0.4pt}}\par\vspace{10pt}
\noindent\begin{minipage}[t]{0.40\textwidth}
\captionof{table}{Jet survey at the body station. Core: within 5\,\% of the centre speed.}
\begin{dbftable}{colspec={Q[c,wd=0.8in] X[c] X[c]}, row{2-Z}={ht=0.16in, valign=m}}
Offset, mm & Horiz., m/s & Vert., m/s \\
$-$60 && \\ $-$45 && \\ $-$30 && \\ $-$15 && \\ 0 && \\ +15 && \\ +30 && \\ +45 && \\ +60 && \\
\end{dbftable}
\end{minipage}\hfill
\begin{minipage}[t]{0.56\textwidth}
\captionof{table}{Test articles. Masses to 0.1~g.}
\begin{dbftable}{colspec={Q[c,wd=0.4in] X[l] Q[c,wd=0.55in] Q[c,wd=0.55in] Q[c,wd=0.55in]}, row{2-Z}={ht=0.28in, valign=m}}
Body & Shape, length, Ø, fin span & Light, g & Mid, g & Heavy, g \\
A &&&& \\ B &&&& \\ C &&&& \\
\end{dbftable}
\end{minipage}
\par\vspace{10pt}
\captionof{table}{Run log. Behaviour: S stable, O oscillating, C coning, R rotating (fail). Roll: none, intermittent, continuous.}
\begin{dbftable}{colspec={Q[c,wd=0.32in] Q[c,wd=0.4in] Q[c,wd=0.6in] Q[c,wd=0.55in] Q[c,wd=0.55in] Q[c,wd=0.5in] Q[c,wd=0.45in] Q[c,wd=0.7in] Q[c,wd=0.55in] Q[c,wd=0.5in] X[l]}, row{2-Z}={ht=0.19in, valign=m}}
Run & Body & Ballast, g & Setting & V, m/s & $\theta$, deg & T, g & Behaviour & Roll & Decay, s & Notes \\
1 &&&&&&&&&& \\ 2 &&&&&&&&&& \\ 3 &&&&&&&&&& \\ 4 &&&&&&&&&& \\ 5 &&&&&&&&&& \\
6 &&&&&&&&&& \\ 7 &&&&&&&&&& \\ 8 &&&&&&&&&& \\ 9 &&&&&&&&&& \\ 10 &&&&&&&&&& \\
11 &&&&&&&&&& \\ 12 &&&&&&&&&& \\ 13 &&&&&&&&&& \\ 14 &&&&&&&&&& \\ 15 &&&&&&&&&& \\
16 &&&&&&&&&& \\
\end{dbftable}
}]

\end{document}
